\begin{statementbox}
\paragraph{用有理分式作局部近似}
记下列五个分式为 $R_E,R_L,R_S,R_C,R_T$。当 $x\to0$ 时，
\[
\begin{aligned}
e^x&=\underbrace{\frac{x^2+6x+12}{x^2-6x+12}}_{R_E(x)}+O(x^5),\\[4pt]
\ln(1+x)&=\underbrace{\frac{3x^2+6x}{x^2+6x+6}}_{R_L(x)}+O(x^5),\\[4pt]
\sin x&=\underbrace{\frac{6x}{x^2+6}}_{R_S(x)}+O(x^5),\\[4pt]
\cos x&=\underbrace{\frac{12-5x^2}{12+x^2}}_{R_C(x)}+O(x^6),\\[4pt]
\tan x&=\underbrace{\frac{3x}{3-x^2}}_{R_T(x)}+O(x^5).
\end{aligned}
\]
$O(x^m)$ 表示误差的绝对值在零点某邻域内不超过 $C|x|^m$，其中 $C$ 为与 $x$ 无关的正常数；它本身不提供具体的有效数字保证。

\textbf{定义域与单位：}指数、正弦、余弦的这些分式对所有实数有定义；对数比较要求 $x>-1$，此时其分母也为正。
正切比较还须 $\cos x\ne0$ 且 $x\ne\pm\sqrt3$，零点附近可取 $|x|<\pi/2$。
所有三角函数的 $x$ 均为\textbf{弧度数}。

\textbf{使用边界：}这些是近似公式，不是恒等式，也不自动给出上界或下界。
本条另证：$0\le x\le1$ 时，
\[
0\le e^x-R_E(x)\le\frac{3x^5}{665},\qquad
0\le\ln(1+x)-R_L(x)\le\frac{x^5}{180}.
\]
两式等号均仅在 $x=0$ 成立。指数在 $x<0$ 时、对数在 $-1<x<0$ 时，误差方向反转。
\end{statementbox}
